\section{The \ecvol{} Benchmark}
\label{sec:benchmark}

We first document the substrate the literature shares (\S\ref{sec:substrate}), then
describe what \ecvol-bench supplies in its place: corpora with reconstructed
identity (\S\ref{sec:corpora}), targets under both horizon conventions
(\S\ref{sec:targets}), measured call timing (\S\ref{sec:timing}), leakage-proof
splits (\S\ref{sec:splits}), and the release (\S\ref{sec:release}).

\subsection{The substrate the literature shares}
\label{sec:substrate}
The EarningsCall (EC) benchmark of \citet{qin2019what} and the two MAEC releases
(2015 and 2016 cohorts; \citealp{li2020maec}) are evaluated, across the lineage in
\S\ref{sec:related}, on split and label files that are redistributed from paper to
paper rather than recomputed. We fetched the canonical files from the VolTAGE, KeFVP,
and HTML repositories, manifested them by SHA-256, and audited four properties.
Table~\ref{tab:substrate} (Result Table 5R) is regenerated from those copies by
\texttt{ecvol audit substrate}; every figure below carries its source file and counts
in the released CSV.
\emph{(1) Provenance.} The VolTAGE and KeFVP test-split files are byte-identical:
the lineage shares one physical label set, never recomputed from prices.
\emph{(2) Embargo.} Every train/validation and validation/test boundary has a
zero-day gap---on EC, training ends and validation starts on the same date
(2017-08-03; 4 training and 14 validation calls on that day)---against targets that
run 30 days past each call, so every boundary leaks the target window.
\emph{(3) Company overlap.} On the temporal splits the field reports, 80.4\% (EC),
44.2\% (MAEC-15), and 55.7\% (MAEC-16) of test calls come from companies that also
appear in training. A temporal split \emph{always} permits this; the finding is not
that overlap exists but that no paper in the lineage reports the disjoint
condition, so the identity confound of \citet{scss2024} has never been measured on
these benchmarks. \emph{(4) Label integrity and convention.} The single-day log-vol
series contain 137 exact-zero cells out of 16{,}800 (0.82\%; the headline averaged
series is clean), and KeFVP's MAEC files interleave four binary labels at exactly
$\tau\in\{3,7,15,30\}$ among 26 price columns under identical
\texttt{future\_label\_N} naming. Finally, the source paper defines its horizons
over \emph{calendar} days; most downstream work, and our own design document until
its correction, describes them as trading days---at $\tau{=}30$ the difference is
roughly 21 versus 30 sessions, and every published number depends on which one was
used. \emph{(5) Labels versus prices.} KeFVP's shipped MAEC test labels key on our
own call identifiers; on the 74--77\% of calls we can price at $\tau\ge7$ they
correlate only 0.51--0.88 with our calendar-day targets (0.25--0.71 at $\tau{=}3$,
where half the windows hold one session), with mean absolute log-vol differences
of 0.15--0.73---so the lineage's labels are not reproducible from prices under
any single documented rule.

% Curated from data/results/result_table_5r.csv and result_table_5r_labels.csv
% (run 20260902T054152Z-audit-substrate-48b9432a). Overlap = test calls whose
% ticker appears in train; embargo = calendar days between the last train (val)
% date and the first val (test) date. Ours = committed splits, same metrics.
\begin{table}[t]
\centering
\small
\begin{tabular}{lrrrrr}
\toprule
Benchmark / split & calls (tr/va/te) & overlap & embargo tr$\to$va & va$\to$te & label defect \\
\midrule
EC (Qin \& Yang)         & 392 / 56 / 112  & 80.4\% & 0 d & 0 d & 137 / 16{,}800 zero cells \\
MAEC-15 (KeFVP files)    & 535 / 76 / 154  & 44.2\% & 0 d & 0 d & 4 binary cols at $\tau$ \\
MAEC-16 (KeFVP files)    & 980 / 140 / 280 & 55.7\% & 0 d & 3 d & 4 binary cols at $\tau$ \\
\midrule
ours: FinCall temporal        & 1480 / 228 / 481 & 91.3\% & 48 d & 48 d & --- \\
ours: FinCall ticker-disjoint & 1701 / 239 / 484 & \textbf{0.0\%} & n/a & n/a & --- \\
ours: MAEC temporal           & 1604 / 73 / 507  & 78.9\% & 101 d & 48 d & --- \\
ours: MAEC ticker-disjoint    & 1805 / 259 / 514 & \textbf{0.0\%} & n/a & n/a & --- \\
ours: Earnings25 ticker-disjoint & 333 / 47 / 95 & \textbf{0.0\%} & n/a & n/a & --- \\
\bottomrule
\end{tabular}
\caption{\textbf{Substrate audit of the shared benchmark files} (Result Table 5R,
regenerated by \texttt{ecvol audit substrate} from SHA-256-manifested copies of the
VolTAGE and KeFVP repositories). The VolTAGE and KeFVP split files are
byte-identical. Overlap is a property every temporal split permits, ours included;
the point is that no paper in the lineage reports the disjoint condition. The
zero-day embargos leak 30-day target windows across every boundary. Shipped
MAEC labels join our calendar-day targets on 74--77\% of test calls at
$\tau\ge7$ and correlate 0.51--0.88 with them (0.25--0.71 at $\tau{=}3$).}
\label{tab:substrate}
\end{table}


\subsection{Corpora and identity reconstruction}
\label{sec:corpora}
We build on two openly licensed multimodal earnings-call corpora and reserve a
third, post-cutoff corpus for the lookahead and clean-audio experiments.
\textbf{FinCall-Surprise} \citep{fincall2025} (primary; Apache-2.0) contains 2{,}688
calls spanning 2019--2021 with full MP3 audio, JSON transcripts, and slides.
\textbf{MAEC} \citep{li2020maec} (secondary; CC-BY-SA-4.0) contains $\sim$3{,}443
calls over 2015--2018 (S\&P~1500) with sentence-aligned audio features and
transcripts but no raw audio. \textbf{Earnings25} \TODO{cite; Zenodo
10.5281/zenodo.18762168} (CC-BY-4.0; $\sim$500 Q4-2025 S\&P~500 calls, 498\,h of
audio, held January--February 2026) is post-cutoff for every model we use and
carries webcast MP3 audio at 64\,kbps for 281 calls and $\le$24\,kbps for 233; it is
the post-cutoff corpus (\S\ref{sec:experiments}). We admit a call when its company
resolves to an SEC ticker that was an S\&P~500 constituent on the release date
(476 of 514; the rest are name-collision look-alikes such as Grainger PLC or
Paramount Group, reason-coded), and 475 join prices.

FinCall-Surprise ships with four fields per record and \emph{no} ticker, company,
CIK, date, time, fiscal period, or speaker names. We reconstruct identity and date
from transcript content and slide metadata, cross-check against price data, and
release the resulting table together with an 80-row manual override list: the
reconstruction resolves 2{,}496/2{,}688 calls (92.9\% overall; 95.3\% of the 2{,}499
earnings-type calls) and passes a 60/60 manual audit. This table is what makes the
identity controls of \S\ref{sec:experiments} runnable on the corpus at all.

\subsection{Price data and targets}
\label{sec:targets}
Daily adjusted-close OHLCV is obtained from a free, no-account source and
cross-checked against Tiingo on a random ticker sample (adjusted-return correlation
$>0.999$); delisted or unmatched tickers produce an explicit coverage report rather
than being silently dropped. After price joining, FinCall-Surprise yields 2{,}424
calls with at least one target (97.1\% join rate) and MAEC yields 2{,}578 (75.4\%; a
documented shortfall driven by 2015--2018 delistings), for a combined modelling set
of 5{,}002 calls.

For each call we define returns $r_t = (P_t - P_{t-1})/P_{t-1}$ relative to day
$t=0$, the last trading day before the call's information is public
(\S\ref{sec:timing}). The primary target is log realised volatility over horizon
$\tau \in \{3,7,15,30\}$,
\begin{equation}
  v_{\mathrm{post}}(\tau) = \ln\!\Big(\sqrt{\tfrac{1}{\tau}\textstyle\sum_{t=1}^{\tau}(r_t-\bar r)^2}\Big),
\end{equation}
with pre-call volatility $v_{\mathrm{pre}}(\tau)$ defined over the sessions
ending at day~0. We ship every target under \emph{both} horizon conventions. Under
the \textbf{trading-day} convention the window is sessions $1..\tau$ ($n=\tau$).
Under the \textbf{calendar-day} convention of \citet{qin2019what}, which the
published leaderboard inherits, the window is every session dated within $\tau$
calendar days of day~0, so $n \approx 0.7\tau$ and varies with the weekday and
holidays (FinCall: $n = 2.4 / 4.9 / 10.7 / 21.0$ sessions on average at
$\tau = 3/7/15/30$). In both cases the variance is taken over the $n$ sessions
actually present, and a window holding fewer than two sessions is excluded with a
reason code rather than reported as a one-return ``volatility''---which removes
\textbf{48.8\%} of FinCall's joined calls at $\tau{=}3$ calendar days (every
Thursday or Friday call). Each convention is a separate file with an identical
schema, a \texttt{convention} column and the per-row session counts; the
trading-day file is the validated set every result in \S\ref{sec:results} uses,
and every result table carries the convention as a column
(\S\ref{sec:results-convention}). We add two identity-robust variants: the \textbf{change} $\Delta v(\tau) = v_{\mathrm{post}}(\tau)
- v_{\mathrm{pre}}(\tau)$, which subtracts the company's own level, and the
\textbf{HAR-residual} $v_{\mathrm{post}}(\tau) - \widehat{\mathrm{HAR}}(\tau)$,
deviation from a train-only HAR-RV \citep{corsi2009har} forecast. Insufficient
post-call history and zero-variance windows yield \texttt{NaN} targets excluded with
a reason code; targets are unit-tested against hand-computed values.

\subsection{Call timing}
\label{sec:timing}
Fewer than 4\% of FinCall-Surprise transcripts mention a clock time, and 32.4\% of
recovered \emph{dates} come from slide-PDF creation stamps, so no call in either
corpus carries a measured time of day. All results in \S\ref{sec:results} therefore
use a uniform, flagged after-hours rule (calls assumed to occur after market close;
$t{=}1$ is the next trading session). We retrofit the information boundary per call from
EDGAR: the acceptance time of the earliest 8-K carrying Item~2.02 (the earnings
release) within three days before to one day after the call date, taken from the
SEC submissions API and cached; where no such filing exists we fall back, in
order, to Earnings25's call-start time, to the before/after-market attribute of
\citet{scss2024}, and finally to the after-hours assumption, with the tier recorded
per call. Coverage by the EDGAR tier is 89.3\% for FinCall-Surprise, 68.6\% for
MAEC, and 99.8\% for Earnings25. Under the measured anchor, \textbf{53\% of FinCall
and 59\% of MAEC calls change day~0}, almost all one session earlier: releases
before the open are the norm (FinCall 53\%, Earnings25 72\%), so the uniform
after-hours rule places the reaction session in the \emph{pre} window for most of
the corpus. Because every released label set in this literature anchors on the
call date, this affects the published targets as much as ours. We keep the
after-hours targets as the primary set for every result in \S\ref{sec:results}
and ship the measured-anchor targets beside them (\texttt{targets\_measured}),
with a sensitivity table (\texttt{timing\_sensitivity.csv}): persistence MSE rises
1--8\% on FinCall and up to 26\% on MAEC at $\tau{=}3$ once the reaction session
enters the post window, while the HAR-RV $R^2_{\mathrm{OOS}}$ moves by at most
$\pm0.08$; the target correlation between anchors is 0.97 for $v_{\mathrm{post}}$
and 0.80 for $\Delta v$ on Earnings25. Whether to re-baseline on the measured
anchor is a release decision we flag rather than take here, since it moves every
split boundary.

\subsection{Leakage-proof splits}
\label{sec:splits}
Three split families are generated once as committed CSVs and guarded by
leakage assertions in CI. \emph{Temporal}: train $<$ validation $<$ test in strict
calendar order ($\sim$70/10/20) with a \textbf{30-trading-day embargo} between
segments so that no 30-day target window crosses a boundary---the constraint the
shared substrate lacks (\S\ref{sec:substrate}). \emph{Ticker-disjoint}: companies in
validation/test never appear in training (grouped, seeded), the headline
identity-control split; by construction its test-call company overlap is 0.0\%,
against 91.3\% (FinCall-Surprise) and 78.9\% (MAEC) on our own temporal splits.
\emph{Combined}: both constraints simultaneously, reported as the hardest
robustness row. An information rule is codified and enforced: every feature function
receives an \texttt{as\_of} timestamp, and the harness fails any feature that reads
price data after it. FinCall-Surprise spans COVID; per-year breakdowns are
mandatory and a Feb--Apr~2020 exclusion is reported as an exploratory sensitivity.

\subsection{Release}
\label{sec:release}
\ecvol-bench is released in full: the identity/date reconstruction table and
override list, split CSVs, targets under both conventions, the derived-feature
archives of \S\ref{sec:methods}, the evaluation harness and baselines, the
substrate-audit report, and every result table. FinCall-Surprise-derived artifacts
are Apache-2.0; MAEC-derived artifacts inherit CC-BY-SA-4.0. Raw audio and
transcripts from sources with unclear licenses are never redistributed (a
\emph{scripts-not-data} policy).
